\chapter{The Z-Transform and the Laplace Transform}

So far the Fourier transform has been our universal lens: feed it a signal, get back
``how much of each frequency is inside.'' That lens is wonderful, but it has a blind
spot. The Fourier transform only ever looks at signals that \emph{neither grow nor
decay} -- pure oscillations of constant amplitude. The instant a signal blows up, or a
filter teeters on the edge of stability, the Fourier integral (or sum) refuses to
converge and the lens goes dark.

This chapter widens the lens. We replace the pure oscillation $e^{\jj\omega t}$ with a
\emph{spiral} $e^{st}$ that is allowed to grow or shrink as it turns. The continuous
version of this idea is the \textbf{Laplace transform}; the discrete version is the
\textbf{z-transform}. Both are built from exactly one new move: let the frequency
variable become \emph{complex}. Everything geometric you already know about the complex
plane -- magnitude as radius, angle as phase, multiplication as rotate-and-scale -- carries
straight over, giving a remarkably compact picture of every linear
time-invariant system: a handful of points on a plane, the \textbf{poles and zeros}.

\begin{focusbox}
The one idea of this chapter: \emph{a frequency that is allowed to grow or decay is just
a complex number}. Plot where a system's response blows up (poles) and where it vanishes
(zeros), and you can read off stability, transient behavior, and the entire frequency
response by walking around a circle. The Fourier transform is the special case where you
only ever stand exactly on that circle.
\end{focusbox}

\section{Why Fourier is not enough}
\label{ch15:why-fourier-not-enough}

Recall the discrete-time Fourier transform (DTFT) of a sequence $x[n]$:
\[
  X(e^{\jj\omega}) = \sum_{n=-\infty}^{\infty} x[n]\, e^{-\jj\omega n}.
\]
The kernel $e^{-\jj\omega n}$ has magnitude exactly $1$ for every $n$. It rotates but
never shrinks. So for this sum to converge, $x[n]$ itself must eventually shrink: the
sequence has to be \emph{absolutely summable}, $\sum_n \abs{x[n]} < \infty$.

That rules out a lot of signals we genuinely care about. The unit step $u[n]$ (all ones
for $n \ge 0$) is not summable. The growing geometric sequence $2^n u[n]$ certainly is
not. An unstable filter's impulse response is not. Yet these are exactly the signals
whose behavior we most need to understand -- often \emph{because} they misbehave.

\begin{intuitionbox}
The DTFT tests your signal against rotating unit-vectors. If your signal grows, no
amount of rotating-without-shrinking can ``catch up'' to it, and the sum diverges. The
remedy emerges naturally once the problem is identified: give the test vectors
permission to grow or shrink too. Multiply the kernel by a real radius $r^{-n}$. If $r$
is big enough, it tames any growth in $x[n]$ before summing. That single extra knob
\emph{is} the z-transform.
\end{intuitionbox}

\section{From rotating to spiraling: defining the z-transform}
\label{ch15:ztransform-defn}

Take the DTFT and replace the unit-magnitude kernel $e^{\jj\omega}$ with a general
complex number $z = r e^{\jj\omega}$, where $r=\abs{z}$ is now free to be anything
positive.

\begin{definition}[The z-transform]
\label{ch15:def-zt}
For a discrete-time signal $x[n]$, the (bilateral) z-transform is
\[
  X(z) \;=\; \Ztr\{x[n]\} \;=\; \sum_{n=-\infty}^{\infty} x[n]\, z^{-n},
  \qquad z \in \C.
\]
The set of complex $z$ for which this sum converges is the \emph{region of convergence}
(ROC).
\end{definition}

Write $z^{-n} = r^{-n} e^{-\jj\omega n}$ and the connection to the Fourier transform is immediate:
\[
  X(re^{\jj\omega}) \;=\; \sum_{n} \bigl( x[n]\, r^{-n}\bigr)\, e^{-\jj\omega n}
  \;=\; \text{DTFT of } \bigl( x[n]\, r^{-n} \bigr).
\]
So the z-transform is nothing but \emph{the Fourier transform of an exponentially
reweighted version of your signal}. Choose the radius $r$ that flattens out your signal's
growth, and the Fourier machinery converges again.

\begin{happeningbox}
Geometrically, the kernel $z^{-n} = (1/r)^{n} e^{-\jj\omega n}$ is a \textbf{spiral}: as
$n$ marches forward it both rotates (the $e^{-\jj\omega n}$ part) and scales (the
$r^{-n}$ part). The DTFT used a circle (constant radius); the z-transform uses a whole
family of spirals, one for every radius. Standing exactly on the unit circle $r=1$
returns you to the DTFT. The z-plane is just the complex plane with this spiral
interpretation pinned to it.
\end{happeningbox}

\begin{connectionbox}
This puts the polar geometry of Chapter~4 directly to work. A complex number $z=re^{\jj\omega}$
packages a radius and an angle together; raising it to the power $-n$ scales the angle to
$-n\omega$ and raises the radius to $r^{-n}$. ``Multiplication is rotate-and-scale'' is why
the z-transform works. If $r$ and $\omega$ still feel like separate pieces, re-read the
polar-form section of Chapter~4 until $re^{\jj\omega}$ feels like a single geometric object.
\end{connectionbox}

\subsection{A worked transform: the right-sided exponential}

The most important example in the whole subject. Let $x[n] = a^{n} u[n]$ -- a geometric
sequence switched on at $n=0$, with $a$ a (possibly complex) constant. Then
\[
  X(z) = \sum_{n=0}^{\infty} a^{n} z^{-n} = \sum_{n=0}^{\infty} \left(\frac{a}{z}\right)^{n}.
\]
Recall that for any common ratio $w$, the standard geometric series $\sum_{n=0}^{\infty} w^n = 1 + w + w^2 + \dots$ converges to $\frac{1}{1-w}$ if and only if $\abs{w} < 1$. Setting $w = a/z = a z^{-1}$, the convergence requirement $\abs{a/z} < 1$ translates directly to $\abs{z} > \abs{a}$. In complex analysis, this summation is a Laurent series in negative powers of $z$; while ordinary Taylor series with positive powers converge inside a circular disk, series in negative powers converge in the exterior region outside a circle. The sum evaluates to
\[
  \boxed{\,X(z) = \frac{1}{1 - a z^{-1}} = \frac{z}{z-a}, \qquad \text{ROC: } \abs{z} > \abs{a}.\,}
\]
Look at what happened. The infinite signal collapsed to a tidy ratio of polynomials in
$z$. There is a value of $z$ that makes the denominator vanish ($z=a$): the function
explodes there. That special point is a \textbf{pole}. And the numerator vanishes at
$z=0$: a \textbf{zero}. The whole signal is now encoded in two dots and an ROC.

\begin{example}[Two transforms, same formula, different ROC]
\label{ch15:ex-roc-matters}
Now try the \emph{left-sided} signal $y[n] = -a^{n}u[-n-1]$ (nonzero only for
$n \le -1$). A short calculation gives
\[
  Y(z) = -\sum_{n=-\infty}^{-1} a^n z^{-n} = \frac{z}{z-a}, \qquad \text{ROC: } \abs{z} < \abs{a}.
\]
The algebraic expression is \emph{identical} to the right-sided case, but the ROC is the
inside of the circle, not the outside. This example underscores the core lesson: \textbf{the formula
$X(z)$ alone does not determine the signal -- you must also state the ROC}. Same poles,
same zeros, opposite signals. The ROC is what tells them apart.
\end{example}

\begin{trapbox}
Beginners memorize ``$a^n u[n] \leftrightarrow z/(z-a)$'' and quietly drop the ROC. Then
they are baffled when an inverse transform is ambiguous or a stability test gives
nonsense. The pole locations are the skeleton; the ROC is which side of each circle you
live on. Always carry the ROC.
\end{trapbox}

\section{The region of convergence, geometrically}
\label{ch15:roc}

The ROC is always a set defined purely by the radius $\abs{z}$ -- never by the angle.
Why? Because convergence of $\sum_n x[n] r^{-n} e^{-\jj\omega n}$ depends on whether
$x[n]r^{-n}$ is summable, and that magnitude test does not involve $\omega$ at all.
Geometrically: \textbf{the ROC is a union of full circles centered at the origin -- it is
always an annulus} (a ring), possibly with inner radius $0$ or outer radius $\infty$.

The shape of the ring follows directly from the ``handedness'' of the signal:

\begin{itemize}
\item \textbf{Right-sided} (causal-ish, $x[n]=0$ for $n<N$): ROC is the \emph{outside} of
a circle, $\abs{z} > r_{\max}$, reaching out to infinity. Growth at large $n$ is what
sets the inner radius.
\item \textbf{Left-sided} (anti-causal): ROC is the \emph{inside} of a circle,
$\abs{z} < r_{\min}$.
\item \textbf{Two-sided}: ROC is a genuine annulus $r_1 < \abs{z} < r_2$ sandwiched
between an inner and outer circle (and may be empty, meaning no z-transform exists).
\item \textbf{Finite-length}: ROC is the entire plane, possibly minus $z=0$ and/or
$z=\infty$.
\end{itemize}

A bedrock rule: \textbf{the ROC never contains a pole}. Poles are exactly where $X(z)$
blows up, so the sum cannot converge there. The poles are therefore the ``fence posts''
between which the ROC ring is stretched.

\begin{figure}[H]
\centering
\begin{tikzpicture}[scale=1.0]
  % ROC annulus (outside pole circle of radius 0.8*1.5=1.2 in drawing units)
  \begin{scope}
    \clip (0,0) circle (2.4);
    \fill[sigblue!12] (0,0) circle (2.4);
    \fill[white] (0,0) circle (1.2);
  \end{scope}
  % axes
  \draw[->,gray!60] (-2.6,0)--(2.6,0) node[right,font=\scriptsize]{Re};
  \draw[->,gray!60] (0,-2.6)--(0,2.6) node[above,font=\scriptsize]{Im};
  % unit circle
  \draw[thick,gray!60] (0,0) circle (1.5);
  \node[gray!70!black,font=\scriptsize,anchor=south west] at (45:1.52) {unit circle};
  % pole circle
  \draw[sigblue,dashed] (0,0) circle (1.2);
  % pole at 0.8 (drawn at radius 1.2)
  \draw[trapred,thick] (1.12,-0.08) -- (1.28,0.08);
  \draw[trapred,thick] (1.12,0.08) -- (1.28,-0.08);
  \node[trapred,font=\scriptsize,anchor=north east] at (1.12,-0.10) {$z=a$};
  % zero at origin
  \draw[sigblue,thick,fill=white] (0,0) circle (0.07);
  \node[sigblue,font=\scriptsize,anchor=south east] at (-0.1,0.1) {zero};
  \node[sigblue!80!black,font=\scriptsize\bfseries] at (1.8,2.1) {ROC: $|z|>|a|$};
\end{tikzpicture}
\caption{Pole-zero plot and ROC for $x[n]=a^n u[n]$ with $\abs{a}=0.8$. A pole
($\times$) sits at $z=a$, a zero ($\circ$) at the origin. The ROC (shaded) is everything
\emph{outside} the pole circle. Because $\abs{a}<1$, the ROC contains the unit circle, so
the DTFT exists and the signal is stable.}
\label{ch15:fig-roc}
\end{figure}

\begin{whybox}
Why is ``ROC contains the unit circle'' the same as ``the DTFT exists''? Because the unit
circle \emph{is} the DTFT, by construction (set $r=1$ in $z=re^{\jj\omega}$). If $r=1$
lies inside the convergence ring, then $\sum_n x[n]e^{-\jj\omega n}$ converges -- which is
precisely the statement that the Fourier transform of $x[n]$ exists. So the single
geometric question ``does the shaded ring cross the unit circle?'' answers ``is this
signal Fourier-analyzable / stable?''
\end{whybox}

\section{Poles, zeros, and the rational picture}
\label{ch15:poles-zeros}

For any system built from delays, adds, and constant multiplies -- which is to say,
\emph{every} linear difference equation -- the z-transform of the impulse response is a
ratio of polynomials:
\[
  H(z) = \frac{B(z)}{A(z)}
       = \frac{b_0 + b_1 z^{-1} + \dots + b_M z^{-M}}{1 + a_1 z^{-1} + \dots + a_N z^{-N}}.
\]
Factor numerator and denominator:
\[
  H(z) = G\,\frac{\prod_{k=1}^{M}(z - z_k)}{\prod_{k=1}^{N}(z - p_k)}.
\]
The roots $z_k$ of the numerator are the \textbf{zeros} (where $H$ vanishes); the roots
$p_k$ of the denominator are the \textbf{poles} (where $H$ blows up); $G$ is an overall
gain. A pole-zero plot draws zeros as $\circ$ and poles as $\times$ on the z-plane. That
small picture, plus the ROC, \emph{is} the system -- it contains everything the difference
equation does.

\begin{focusbox}
\textbf{Poles control what the system does on its own} (the natural/transient response,
resonances, stability). \textbf{Zeros control what the system rejects} (notches, nulls,
the shaping of the response). Read poles for ``how it rings,'' zeros for ``what it
kills.''
\end{focusbox}

\subsection{Difference equation $\to$ transfer function, mechanically}

The z-transform turns the awkward operation of \emph{time delay} into plain
multiplication, via the delay property:
\[
  x[n-k] \;\xleftrightarrow{\ \Ztr\ }\; z^{-k}\,X(z).
\]
So $z^{-1}$ is literally ``the delay-by-one-sample operator.'' This single fact converts
any difference equation into algebra.

\begin{example}[A one-pole lowpass]
\label{ch15:ex-onepole}
Consider the smoothing filter $y[n] = \alpha\,y[n-1] + (1-\alpha)\,x[n]$ with
$0<\alpha<1$ (an exponential moving average). Transform both sides using the delay
property:
\[
  Y(z) = \alpha z^{-1} Y(z) + (1-\alpha) X(z)
  \ \Longrightarrow\
  H(z) = \frac{Y(z)}{X(z)} = \frac{1-\alpha}{1 - \alpha z^{-1}} = \frac{(1-\alpha)\,z}{z - \alpha}.
\]
One pole at $z=\alpha$, one zero at $z=0$. Since $0<\alpha<1$, the pole is inside the unit
circle, so a causal (right-sided) implementation has ROC $\abs{z}>\alpha$ which contains
the unit circle: the filter is stable. The closer $\alpha$ creeps toward $1$, the closer
the pole creeps to the unit circle at $z=1$ (DC), and the more sharply the filter
emphasizes low frequencies -- and the longer it rings. We will see this ``ringing'' is
exactly the pole pulling on the frequency response.
\end{example}

\begin{connectionbox}
Notice the impulse response of that filter is $h[n]=(1-\alpha)\alpha^n u[n]$ -- a decaying
exponential -- and its transform is the right-sided $z/(z-\alpha)$ form from
Section~\ref{ch15:ztransform-defn}. \emph{The pole location $\alpha$ is the base of the
exponential.} Pole near the origin $\Rightarrow$ fast decay; pole near the unit circle
$\Rightarrow$ slow decay; pole on or outside it $\Rightarrow$ no decay or blow-up. The
pole's radius is the system's memory.
\end{connectionbox}

\begin{figure}[H]
\centering
\begin{tikzpicture}
\begin{axis}[
  width=11cm, height=5cm,
  title={Impulse response for two pole radii},
  title style={font=\small},
  xlabel={$n$ (samples)}, ylabel={$h[n]$},
  xmin=-0.5, xmax=24, ymin=0, ymax=0.55,
  ytick={0,0.2,0.4},
  axis lines=left,
  legend pos=north east, legend cell align=left,
]
% h[n] = (1-alpha) alpha^n u[n]
\addplot[ycomb,sigblue,thick,mark=*,mark size=1.2pt,domain=0:23,samples=24]
  {0.5*pow(0.5,x)};
\addlegendentry{$\alpha=0.5$ (pole at $0.5$)}
\addplot[ycomb,trapred,thick,mark=*,mark size=1.2pt,domain=0:23,samples=24]
  {0.1*pow(0.9,x)};
\addlegendentry{$\alpha=0.9$ (pole at $0.9$)}
\end{axis}
\end{tikzpicture}
\caption{Impulse response $h[n]=(1-\alpha)\alpha^n u[n]$ of the one-pole lowpass for two
pole radii. The pole at $\alpha=0.5$ (blue) sits closer to the origin and decays fast; the
pole at $\alpha=0.9$ (red) hugs the unit circle and decays slowly, so the filter holds a
longer memory and rings longer. The pole's radius is exactly the per-sample decay factor.}
\label{ch15:fig-impulse}
\end{figure}

\section{Stability and causality, read off the plane}
\label{ch15:stability}

This is where the pole-zero plot pays off directly. Two questions dominate filter design:
\emph{is it stable?} and \emph{is it causal (can I run it forward in real time)?} Both
are answered by where the poles sit relative to the unit circle, together with the ROC.

\begin{formalbox}
\textbf{Stability (BIBO).} A causal LTI system is stable if and only if all its poles lie
\emph{strictly inside} the unit circle, $\abs{p_k}<1$. Equivalently: the ROC includes the
unit circle. \\[0.3em]
\textbf{Causality.} A system is causal if and only if its ROC is the exterior of the
outermost pole (and $H(z)$ is proper). \\[0.3em]
\textbf{Both at once.} A causal, stable system has all poles inside the unit circle, and
the ROC is the region outside the outermost pole, which therefore necessarily contains the
unit circle.
\end{formalbox}

\begin{intuitionbox}
Why ``inside the unit circle = stable''? A pole at radius $\abs{p}$ contributes a term
$p^n$ to the impulse response. If $\abs{p}<1$, $p^n\to 0$: the natural response dies out,
energy leaks away, the system settles. If $\abs{p}>1$, $p^n$ grows without bound: the
slightest input rings forever louder -- unstable. If $\abs{p}=1$, $p^n$ neither grows nor
dies (a pure undamped oscillation, like an ideal resonator or integrator): marginally
stable, living right on the knife's edge. The unit circle is literally the boundary
between decay and growth.
\end{intuitionbox}

\begin{figure}[H]
\centering
\begin{tikzpicture}[scale=2.2]
  \draw[->,gray!60] (-1.45,0)--(1.45,0) node[right,font=\scriptsize]{Re};
  \draw[->,gray!60] (0,-1.35)--(0,1.35) node[above,font=\scriptsize]{Im};
  \draw[thick,gray!60] (0,0) circle (1);
  \node[gray!70!black,font=\scriptsize,anchor=south west] at (50:1.05) {$|z|=1$};
  % stable pole pair (inside)
  \node[targetgreen] at (0.55,0.45) {\large $\times$};
  \node[targetgreen] at (0.55,-0.45) {\large $\times$};
  \node[targetgreen,font=\scriptsize,right] at (0.6,0.5) {stable};
  % marginal pole
  \draw[quotientorange,thick] ({cos(150)-0.045},{sin(150)-0.045}) -- ({cos(150)+0.045},{sin(150)+0.045});
  \draw[quotientorange,thick] ({cos(150)-0.045},{sin(150)+0.045}) -- ({cos(150)+0.045},{sin(150)-0.045});
  \node[quotientorange,font=\scriptsize,anchor=south east] at ({cos(150)-0.05},{sin(150)+0.05}) {marginal};
  % unstable pole
  \node[trapred] at (-0.65,-0.95) {\large $\times$};
  \node[trapred,font=\scriptsize,below] at (-0.65,-1.02) {unstable};
\end{tikzpicture}
\caption{Stability at a glance. Poles inside the unit circle (green) give a decaying,
stable response. A pole exactly on the circle (orange) is marginally stable -- an
undamped oscillation. A pole outside (red) means the response grows without bound:
unstable. For a \emph{causal} system you read this directly; zeros may sit anywhere
without affecting stability.}
\label{ch15:fig-stability}
\end{figure}

\begin{trapbox}
Zeros do \emph{not} affect stability. You can scatter zeros anywhere -- inside, on, or
outside the unit circle -- and a causal system stays stable as long as the poles behave.
(Zeros outside the circle do affect \emph{phase}: such a filter is called ``non-minimum-phase''
because another filter with the reflected inside-the-circle zero has the same magnitude response with smaller phase lag.) Only poles vote on stability.
\end{trapbox}

\section{The frequency response is a walk around the unit circle}
\label{ch15:freq-response}

Here is the most beautiful payoff of the whole chapter, and the reason DSP engineers
sketch pole-zero plots on napkins. The frequency response is just the z-transform
evaluated on the unit circle:
\[
  H(e^{\jj\omega}) = H(z)\big|_{z=e^{\jj\omega}}
  = G\,\frac{\prod_k (e^{\jj\omega}-z_k)}{\prod_k (e^{\jj\omega}-p_k)}.
\]
Each factor like $(e^{\jj\omega}-p_k)$ is a \emph{vector in the complex plane}: it points
from the pole $p_k$ to the point $e^{\jj\omega}$ on the unit circle. Its magnitude is the
\emph{distance} from the pole to that point. So:
\[
  \abs{H(e^{\jj\omega})}
  = \abs{G}\,\frac{\prod_k \bigl(\text{distance from } e^{\jj\omega} \text{ to zero } z_k\bigr)}
                  {\prod_k \bigl(\text{distance from } e^{\jj\omega} \text{ to pole } p_k\bigr)}.
\]

\begin{happeningbox}
Imagine walking a pin once around the unit circle as $\omega$ goes from $0$ to $2\pi$. At
each step, stretch a rubber band from the pin to every pole and every zero. The gain is
(product of zero-distances) divided by (product of pole-distances). \\[0.3em]
\textbf{When you pass near a pole}, that pole-distance gets tiny, you are dividing by a
small number, and the gain \emph{spikes} -- a resonance/peak. \\[0.3em]
\textbf{When you pass near a zero}, that zero-distance gets tiny, you multiply by a small
number, and the gain \emph{dips} -- a notch (an exact null if the zero is right on the
circle). \\[0.3em]
The entire frequency response is governed by how close your walk around the unit circle
passes to each pole and zero.
\end{happeningbox}

\begin{figure}[H]
\centering
\begin{tikzpicture}[scale=2.2]
  \draw[->,gray] (-1.5,0)--(1.55,0) node[right,font=\scriptsize]{Re};
  \draw[->,gray] (0,-1.35)--(0,1.45) node[above,font=\scriptsize]{Im};
  \draw[thick,gray!60] (0,0) circle (1);
  % pole near circle (resonance)
  \coordinate (P) at (0.78,0.55);
  \node[trapred] at (P) {\large $\times$};
  \node[trapred,font=\scriptsize,right] at (0.82,0.6) {pole $p$};
  % zero on circle
  \coordinate (Z) at ({cos(-40)},{sin(-40)});
  \draw[sigblue,thick] (Z) circle (0.045);
  \node[sigblue,font=\scriptsize,below right] at (Z) {zero $z$};
  % evaluation point on circle
  \coordinate (E) at ({cos(45)},{sin(45)});
  \fill[black] (E) circle (0.025);
  \node[font=\scriptsize,above right] at (E) {$e^{\jj\omega}$};
  % rubber bands
  \draw[trapred,thick,-{Stealth}] (P)--(E);
  \draw[sigblue,thick,-{Stealth}] (Z)--(E);
  \node[trapred,font=\scriptsize] at (1.0,1.1) {small $\Rightarrow$ peak};
  \node[sigblue,font=\scriptsize] at (-0.55,0.95) {large $\Rightarrow$ no dip here};
\end{tikzpicture}
\caption{The frequency response as distances. As $\omega$ sweeps, the evaluation point
$e^{\jj\omega}$ slides around the unit circle. The gain at each $\omega$ is the product of
distances to the zeros divided by the product of distances to the poles. Near the pole $p$
the pole-distance shrinks and the gain peaks; near the zero $z$ (here on the circle) the
gain would null out.}
\label{ch15:fig-rubber}
\end{figure}

\begin{example}[A resonator from a single pole pair]
\label{ch15:ex-resonator}
Put a pair of complex-conjugate poles at $p = r e^{\pm\jj\omega_0}$ with $r$ just under
$1$, say $r=0.95$, $\omega_0 = \pi/4$. The transfer function is
\[
  H(z) = \frac{1}{(1 - re^{\jj\omega_0}z^{-1})(1 - re^{-\jj\omega_0}z^{-1})}
       = \frac{1}{1 - 2r\cos\omega_0\, z^{-1} + r^2 z^{-2}}.
\]
The frequency response peaks sharply at $\omega=\omega_0$, because that is where the
circular walk passes closest to the poles. The peak's sharpness (its Q) is set by
$1-r$: poles hugging the circle give a tall, narrow resonance; poles further in give a
broad, gentle hump. We have designed a tunable bandpass filter from two dots.
\end{example}

\begin{figure}[H]
\centering
\begin{tikzpicture}[scale=1.7]
  \draw[->,gray!60] (-1.35,0)--(1.45,0) node[right,font=\scriptsize]{Re};
  \draw[->,gray!60] (0,-1.3)--(0,1.35) node[above,font=\scriptsize]{Im};
  \draw[thick,gray!60] (0,0) circle (1);
  \node[gray!70!black,font=\scriptsize,anchor=south east] at (135:1.05) {$|z|=1$};
  % conjugate poles at r=0.8, w0=45deg
  \coordinate (P1) at ({0.8*cos(45)},{0.8*sin(45)});
  \coordinate (P2) at ({0.8*cos(45)},{-0.8*sin(45)});
  \coordinate (E) at ({cos(45)},{sin(45)});
  \node[trapred] at (P1) {\large $\times$};
  \node[trapred] at (P2) {\large $\times$};
  \node[trapred,font=\scriptsize,below left] at (P1) {$p_1$};
  \node[trapred,font=\scriptsize,above left] at (P2) {$p_2$};
  % evaluation point & distance vector
  \fill[sigblue] (E) circle (0.035);
  \draw[sigblue,thick,-{Stealth}] (P1)--(E);
  \node[sigblue,font=\scriptsize,above right] at (E) {$e^{\jj\omega_0}$};
  % double zero at origin
  \draw[sigblue,thick] (0,0) circle (0.05);
  \node[sigblue,font=\scriptsize,anchor=north east] at (-0.08,-0.08) {zeros at $0$};
  % angle arc
  \draw[gray,->] (0.45,0) arc (0:45:0.45);
  \node[gray,font=\scriptsize] at (0.6,0.18) {$\omega_0$};
\end{tikzpicture}
\hfill
\begin{tikzpicture}
\begin{axis}[
  width=8.5cm, height=5.2cm,
  title={Magnitude response $|H(e^{\jj\omega})|$},
  title style={font=\small},
  xlabel={$\omega$ (rad/sample)}, ylabel={$|H(e^{\jj\omega})|$},
  domain=0:3.14159, samples=180,
  xmin=0, xmax=3.14159, ymin=0, ymax=16,
  xtick={0,0.7853,1.5708,2.3562,3.1416},
  xticklabels={$0$,$\omega_0{=}\frac{\pi}{4}$,$\frac{\pi}{2}$,$\frac{3\pi}{4}$,$\pi$},
  axis lines=left, every axis plot/.append style={thick},
]
\addplot[sigblue] {1/sqrt( (1 - 2*0.95*0.70711*cos(deg(x)) + 0.9025*cos(deg(2*x)))^2 + (2*0.95*0.70711*sin(deg(x)) - 0.9025*sin(deg(2*x)))^2 )};
\addplot[trapred,dashed,thick] coordinates {(0.7853,0)(0.7853,16)};
\node[trapred,font=\scriptsize,above right] at (axis cs:0.7853,13.5) {resonant peak};
\end{axis}
\end{tikzpicture}
\caption{Geometric alignment: pole-zero z-plane (left, pole gap exaggerated for clarity) aligned with the frequency response magnitude for $r=0.95$, $\omega_0=\pi/4$ (right). When the evaluation point $e^{\jj\omega}$ reaches $\omega_0$, the distance $|e^{\jj\omega}-p_1| = 1-r = 0.05$ shrinks to its minimum, producing the sharp resonant peak $|H(e^{\jj\omega_0})| = 1/\bigl[(1-r)\,|e^{\jj\omega_0}-p_2|\bigr] \approx 14.5$ in the frequency domain.}
\label{ch15:fig-resonator-pz}
\label{ch15:fig-resonator-mag}
\end{figure}

\begin{whybox}
Why complex-\emph{conjugate} pairs? A real difference equation has real coefficients, and
the roots of a real polynomial come in conjugate pairs. Geometrically, a conjugate pair
sits symmetrically above and below the real axis, so their two resonances land at
$+\omega_0$ and $-\omega_0$ -- which is exactly the symmetry a real signal's spectrum
already has. A lone complex pole would describe a complex-valued filter; pairing it up
keeps everything real. This is Chapter~4's ``conjugates undo imaginary
parts'' showing up as hardware reality.
\end{whybox}

\section{The Laplace transform: the continuous twin}
\label{ch15:laplace}

Everything above has a continuous-time mirror image. For continuous signals we cannot use
$z^{-n}$ (there is no integer $n$), so we use the continuous spiral $e^{-st}$ where
$s = \sigma + \jj\omega$ is a complex frequency: $\sigma$ is a growth/decay rate (nepers
per second) and $\omega$ is the ordinary angular frequency.

\begin{definition}[The Laplace transform]
\label{ch15:def-laplace}
For a continuous-time signal $x(t)$, the (bilateral) Laplace transform is
\[
  X(s) = \Ltr\{x(t)\} = \int_{-\infty}^{\infty} x(t)\, e^{-st}\,\dd t, \qquad s=\sigma+\jj\omega \in \C,
\]
with an ROC consisting of the $\sigma$ for which the integral converges.
\end{definition}

Set $\sigma=0$ (i.e. $s=\jj\omega$, evaluate on the imaginary axis) and you recover the
continuous-time Fourier transform $X(\jj\omega)=\int x(t)e^{-\jj\omega t}\dd t$. The
$e^{-\sigma t}$ factor is the continuous analog of the $r^{-n}$ taming factor: it weights
the signal so the integral can converge even when $x(t)$ grows.

\begin{focusbox}
The dictionary between the two worlds:
\begin{center}
\begin{tabular}{lll}
\toprule
 & \textbf{Continuous (Laplace)} & \textbf{Discrete (z)} \\
\midrule
complex frequency & $s=\sigma+\jj\omega$ & $z=re^{\jj\omega}$ \\
Fourier lives on & imaginary axis $\sigma=0$ & unit circle $r=1$ \\
stable region & left half-plane $\mathrm{Re}(s)<0$ & inside circle $\abs{z}<1$ \\
``do nothing'' point & $s=0$ & $z=1$ \\
derivative / delay & $\tfrac{\dd}{\dd t}\leftrightarrow s$ & delay $\leftrightarrow z^{-1}$ \\
\bottomrule
\end{tabular}
\end{center}
\end{focusbox}

\begin{happeningbox}
The single most useful fact about Laplace: \emph{differentiation becomes multiplication by
$s$}. If $x(t)\leftrightarrow X(s)$ then $x'(t)\leftrightarrow sX(s)$ (zero initial
conditions). That is why the Laplace transform simplifies differential equations so effectively -- a circuit's
integro-differential equations turn into ordinary algebra in $s$, you solve, then
transform back. The z-transform does the same trick for difference equations with
$z^{-1}$ as the delay.
\end{happeningbox}

\subsection{Stability in the s-plane}

The mapping of stability is just as clean, but the boundary changes shape. In the
z-plane the dividing line is the unit circle; in the s-plane it is the \textbf{imaginary
axis}. A pole at $s=\sigma+\jj\omega$ contributes $e^{\sigma t}e^{\jj\omega t}$ to the
response: if $\sigma<0$ it decays (stable), if $\sigma>0$ it grows (unstable), if
$\sigma=0$ it oscillates forever (marginal). So for a causal continuous system:
\[
  \text{stable} \iff \text{all poles in the left half-plane } \mathrm{Re}(s_k)<0.
\]

\begin{figure}[H]
\centering
\begin{tikzpicture}[scale=1.0]
  % shaded left half plane
  \fill[targetgreen!12] (-2.7,-1.8) rectangle (0,1.8);
  \draw[->,gray] (-2.8,0)--(2.4,0) node[right,font=\scriptsize]{$\sigma=\mathrm{Re}(s)$};
  \draw[->,gray] (0,-1.9)--(0,1.9) node[above,font=\scriptsize]{$\jj\omega$};
  \draw[thick,gray!70] (0,-1.85)--(0,1.85);
  \node[targetgreen!60!black,font=\scriptsize] at (-1.7,1.6) {stable (LHP)};
  \node[trapred!70!black,font=\scriptsize] at (1.3,1.6) {unstable (RHP)};
  % stable conjugate poles
  \node[targetgreen] at (-1.0,0.8) {\large $\times$};
  \node[targetgreen] at (-1.0,-0.8) {\large $\times$};
  % unstable pole
  \node[trapred] at (0.9,-0.5) {\large $\times$};
  % marginal pole on axis
  \node[quotientorange] at (0,1.2) {\large $\times$};
  \node[quotientorange,font=\scriptsize,right] at (0.07,1.2) {marginal};
\end{tikzpicture}
\caption{The s-plane stability picture. Poles in the left half-plane (green region) give
decaying, stable behavior; poles in the right half-plane grow; poles on the imaginary
axis oscillate undamped. Compare Figure~\ref{ch15:fig-stability}: the unit circle has
been ``unrolled'' into the imaginary axis, its interior into the left half-plane.}
\label{ch15:fig-splane}
\end{figure}

\begin{whybox}
Why does the inside of the circle map to the left half-plane? Because sampling links the
two by $z = e^{sT}$, where $T$ is the sample period. Write $s=\sigma+\jj\omega$:
\[
  z = e^{\sigma T} e^{\jj\omega T} \;\Rightarrow\; \abs{z}=e^{\sigma T}.
\]
So $\sigma<0$ (left half-plane) gives $\abs{z}<1$ (inside the circle), $\sigma=0$ (the
imaginary axis) gives $\abs{z}=1$ (the unit circle), and $\sigma>0$ gives $\abs{z}>1$. The
exponential $e^{sT}$ literally wraps the s-plane's vertical strips around the z-plane's
circles. This is the seed of digital filter design: take an analog filter you trust in the
s-plane and map it into the z-plane.
\end{whybox}

\begin{historybox}
Pierre-Simon Laplace introduced his integral transform in the 1780s while investigating
probability distributions and celestial mechanics. But when radar tracking, guided missiles,
and sampled-data control systems emerged in the 1950s, engineers discovered that continuous-time
calculus could not accurately model control loops where measurements arrive only at discrete clock ticks.
At Columbia University, John Ragazzini, Lotfi Zadeh, and Eliahu Jury pioneered the $z$-transform to
bridge this gap, mapping continuous differential systems into discrete difference equations.

A persistent conceptual trap was the Region of Convergence. In early applications, practitioners
frequently treated the algebraic rational formula $X(z)$ as if it were the entire story. They soon
encountered baffling paradoxes: the causal sequence $a^n u[n]$ and the anti-causal sequence
$-a^n u[-n-1]$ yield the exact same algebraic formula $1/(1 - a z^{-1})$. Without specifying
whether the region of convergence lies outside or inside the boundary circle $|z|=|a|$, one cannot
tell whether the signal is a causal wave decaying into the future or an anti-causal wave exploding
into the past.

Recognizing that $X(z)$ and its region of convergence are inseparable unlocked a visual revolution:
pole-zero geometry on the complex plane. By mapping the continuous frequency axis onto the discrete
unit circle $|z|=1$, stability became a simple question of confinement inside the circle, while the
frequency response was revealed as the ratio of geometric distances from poles and zeros to $e^{\jj\omega}$.
Chapter~16 (\emph{Designing Filters}) translates this geometry into standard Butterworth, Chebyshev,
and optimal FIR designs, while Chapter~17 (\emph{Practical Filtering}) addresses the finite-precision
arithmetic needed to run them safely in hardware.
\end{historybox}

\section{Transfer functions and block diagrams}
\label{ch15:transfer-functions}

A \textbf{transfer function} $H(z)$ (or $H(s)$) is the z- (or Laplace-) transform of a
system's impulse response. Equivalently, it is the ratio output/input in the transform
domain: $Y(z)=H(z)X(z)$. The convolution in time becomes multiplication in the transform
domain -- this is the convolution theorem from Chapter~11, now extended off the
circle.

This multiplicative structure is why transfer functions are the lingua franca of system
design. Cascade two systems and their transfer functions \emph{multiply}: $H(z)=H_1(z)
H_2(z)$. Put them in parallel and they \emph{add}: $H(z)=H_1(z)+H_2(z)$. Wrap one in a
feedback loop with $H_2$ and you get the feedback formula $H = H_1/(1+H_1 H_2)$. Complex
systems are built like LEGO from these three rules.

\begin{figure}[H]
\centering
\begin{tikzpicture}[scale=1.0,every node/.style={font=\small}]
  % nodes
  \node (x) at (0,0) {$X(z)$};
  \node[draw,rounded corners,fill=sigblue!10,minimum height=0.8cm] (h1) at (2.2,0) {$H_1(z)$};
  \node[draw,rounded corners,fill=sigblue!10,minimum height=0.8cm] (h2) at (5.0,0) {$H_2(z)$};
  \node (y) at (7.4,0) {$Y(z)$};
  \draw[-{Stealth}] (x)--(h1);
  \draw[-{Stealth}] (h1)--(h2) node[midway,above,font=\scriptsize]{$X H_1$};
  \draw[-{Stealth}] (h2)--(y);
  \node[font=\scriptsize] at (3.6,-0.9) {cascade: $H(z)=H_1(z)\,H_2(z)$};
\end{tikzpicture}
\caption{A cascade of two systems. In the transform domain their transfer functions
simply multiply, and their pole-zero plots overlay (the combined system has the union of
all poles and all zeros). This is why filter design so often proceeds one second-order
``biquad'' section at a time.}
\label{ch15:fig-cascade}
\end{figure}

\begin{applicationbox}
\textbf{Analog circuits.} An RC lowpass has transfer function $H(s)=1/(1+sRC)$: one pole
at $s=-1/RC$ on the negative real axis, deep in the stable left half-plane. The pole's
distance from the origin sets the cutoff frequency $\omega_c=1/RC$. An RLC circuit adds a
conjugate pole pair and can resonate -- exactly the continuous cousin of
Example~\ref{ch15:ex-resonator}. Every analog filter you will ever meet is a pattern of
dots in the s-plane.
\end{applicationbox}

\section{Inverting the transform: partial fractions}
\label{ch15:inversion}

To get back to the time domain you usually do not compute a contour integral -- you do
\textbf{partial fractions}, breaking $H(z)$ into a sum of one-pole pieces whose inverses
you already know.

\begin{example}[Inverse by partial fractions]
\label{ch15:ex-pf}
Take
\[
  H(z) = \frac{1}{(1-\tfrac12 z^{-1})(1-\tfrac13 z^{-1})}.
\]
Expand into partial fractions in $z^{-1}$:
\[
  H(z) = \frac{A}{1-\tfrac12 z^{-1}} + \frac{B}{1-\tfrac13 z^{-1}},
  \quad A = 3,\ B=-2.
\]
Assuming a causal (right-sided) system, each term inverts via the standard pair
$\frac{1}{1-az^{-1}}\leftrightarrow a^n u[n]$:
\[
  h[n] = 3\left(\tfrac12\right)^n u[n] - 2\left(\tfrac13\right)^n u[n].
\]
Both poles ($\tfrac12,\tfrac13$) are inside the unit circle, so $h[n]\to 0$: stable, as
the pole positions promised. The decay is dominated by the pole nearest the unit circle
($\tfrac12$), the slowest-decaying mode.
\end{example}

\begin{trapbox}
The same $H(z)$ with a different ROC inverts to a different $h[n]$ (left-sided terms use
$-a^n u[-n-1]$). Always pin the ROC before inverting, and choose the term form
(right- vs left-sided) to match it. Forgetting this is the number-one inverse-transform
mistake.
\end{trapbox}

\section{Toward the wrist tracker}
\label{ch15:wrist}

\begin{goalbox}
Every filter on the wrist controller -- the lowpass that removes hand tremor from the
accelerometer, the highpass that strips gravity to leave dynamic motion, the bandpass that
isolates a tap's frequency band, the smoother on the EMG envelope -- is a little
pole-zero plot living in a microcontroller. When you write
\begin{center}\ttfamily y += alpha*(x - y);\end{center}
on the Arduino, you have placed a single pole at $z=1-\alpha$ (an exponential smoother,
Example~\ref{ch15:ex-onepole}). Choosing $\alpha$ \emph{is} choosing where that pole sits
between the origin (snappy, noisy) and the unit circle at $z=1$ (smooth, sluggish, laggy).
The pole-zero picture turns ``how laggy should my gesture detector feel?'' into a single
geometric dial.

The stability rule matters in hardware specifically. An IIR (recursive) filter with a
pole accidentally pushed outside the unit circle by a coefficient typo, or by fixed-point
rounding on an 8-bit MCU, will \emph{self-oscillate} -- your gesture signal explodes into
garbage with no input at all. ``Keep the poles inside the circle'' is not abstract; it is
the difference between a filter that works and a sensor that screams.
\end{goalbox}

\begin{buildbox}
A practical second-order ``biquad'' bandpass to isolate a finger-tap transient
($\omega_0$ tuned to the tap's dominant band, $r$ near $0.9$--$0.97$ for selectivity),
realized directly from the two-pole resonator of Example~\ref{ch15:ex-resonator}:
\begin{lstlisting}[language=C]
// One biquad bandpass section, Direct Form I, float math.
// Poles at r*e^{+-j w0}; a zero pair shapes the passband.
float b0, b1, b2, a1, a2;   // precomputed from r, w0
float x1=0, x2=0, y1=0, y2=0;

float biquad(float x) {
    float y = b0*x + b1*x1 + b2*x2 - a1*y1 - a2*y2;
    x2 = x1; x1 = x;        // shift input history
    y2 = y1; y1 = y;        // shift output history
    return y;
}
// Coeffs for a resonator bandpass:
//   a1 = -2*r*cos(w0);  a2 = r*r;        // denominator -> the two poles
//   b0 = 1 - r;  b1 = 0;  b2 = -(1-r);   // zeros at z=+-1: DC/Nyquist nulls
\end{lstlisting}
Sanity check on-device: with $r<1$ the recursion must settle to zero after an impulse. If
it rings forever or grows, a coefficient is wrong and a pole has escaped the circle. On
fixed-point hardware, leave headroom: keep $r \le 0.97$ so rounding cannot nudge
$\abs{p}$ past $1$.
\end{buildbox}

\begin{trapbox}
\textbf{The hard sensor limit.} None of this transform machinery conjures information the
sensor never captured. A wrist-worn IMU senses the forearm's gross motion and orientation;
no amount of clever pole placement recovers \emph{individual finger joint angles} from it
-- that signal simply is not in the accelerometer stream. Distinguishing which finger
tapped needs forearm \emph{EMG} or \emph{force-myography} (muscle/tendon signals), and
\emph{continuous} finger pose needs sensors on the fingers themselves (e.g. instrumented
rings). The z-transform is how you clean and shape whatever the sensor genuinely sees; it
is not a substitute for the right sensor.
\end{trapbox}

\section{Summary}
\label{ch15:summary}

\begin{itemize}
\item The z-transform $X(z)=\sum_n x[n]z^{-n}$ generalizes the DTFT by letting the
frequency variable be a complex number $z=re^{\jj\omega}$ -- a \emph{spiral} kernel that
can grow or shrink. It is the DTFT of the reweighted signal $x[n]r^{-n}$.
\item The \textbf{region of convergence} is always an annulus (ring) and never contains a
pole. Right-sided signals give outside-the-circle ROCs, left-sided give inside. The
formula plus the ROC together determine the signal.
\item A system's transfer function is a ratio of polynomials; its \textbf{poles}
($\times$) and \textbf{zeros} ($\circ$) on the z-plane, plus the ROC, encode it
completely. Poles set resonance and stability; zeros set notches.
\item \textbf{Stability + causality} $\iff$ all poles strictly inside the unit circle (ROC
contains the unit circle). The unit circle is the boundary between decay and growth.
\item The \textbf{frequency response} is the z-transform on the unit circle; its magnitude
is the product of zero-distances over pole-distances as you walk the circle -- peaks near
poles, dips near zeros.
\item The \textbf{Laplace transform} is the continuous analog ($e^{-st}$ kernel,
$s=\sigma+\jj\omega$). Fourier lives on the imaginary axis; stability is the left
half-plane. Sampling links the two via $z=e^{sT}$, wrapping the LHP onto the disk.
\end{itemize}

\begin{connectionbox}
\textbf{Capabilities, Downsides, and Forward Roadmap}
\begin{itemize}
  \item \textbf{What this unlocked}: The $z$-transform and Laplace transform generalize Fourier analysis
  to growing, decaying, and transient signals. By treating frequency as a complex variable $z = r e^{\jj\omega}$,
  system behavior maps to pole-zero geometry: for causal systems, stability means all poles stay strictly
  inside the unit circle, and the frequency response is evaluated by walking the circle.
  \item \textbf{Honest downsides and limitations}: Knowing where poles and zeros \emph{should} sit does not
  instantly produce an optimal filter. Placing poles by hand for sharp roll-offs without causing ripple,
  instability, or excessive phase distortion requires systematic approximation theory.
  \item \textbf{Where we go next}: Chapter~16 translates this geometric framework into \textbf{Filter Design}.
  We explore classical analog prototypes (Butterworth, Chebyshev, elliptic), map them into discrete-time IIR
  filters via the bilinear transform, and design linear-phase FIR filters using windowing and the Parks--McClellan
  algorithm. Chapter~17 then addresses real-time implementation challenges on embedded hardware.
\end{itemize}
\end{connectionbox}

\section{Exercises}
\label{ch15:exercises}

\begin{exercise}
Compute the z-transform and ROC of $x[n] = \left(\tfrac34\right)^n u[n]$. Mark its pole
and zero on a sketch of the z-plane and state whether the DTFT exists.
\end{exercise}

\begin{exercise}
The sequence $x[n] = \left(\tfrac34\right)^{n}u[n] + 2^{n}u[n]$ is a sum of a decaying and
a growing term. Find $X(z)$, locate both poles, and determine the ROC. Does the DTFT
exist? Explain geometrically using the unit circle.
\end{exercise}

\begin{exercise}
Two systems have the \emph{same} pole-zero diagram (a single pole at $z=0.6$, a zero at
the origin) but different ROCs. Write the two possible impulse responses and say which one
is stable and which is causal.
\end{exercise}

\begin{exercise}
A difference equation reads $y[n] = 0.9\,y[n-1] - 0.81\,y[n-2] + x[n]$. Find $H(z)$, its
two poles (in polar form $re^{\pm\jj\omega_0}$), and predict the frequency $\omega_0$ at
which the magnitude response peaks. Is the system stable?
\end{exercise}

\begin{exercise}
For the one-pole lowpass $H(z)=(1-\alpha)/(1-\alpha z^{-1})$, compute
$\abs{H(e^{\jj\omega})}$ at $\omega=0$ and $\omega=\pi$, and show the ratio (the
DC-to-Nyquist attenuation) grows as $\alpha\to 1$. Interpret using the pole-to-circle
distance.
\end{exercise}

\begin{exercise}
Using partial fractions, find the causal inverse z-transform of the rational transfer function
$H(z)=1/\bigl[(1-\tfrac14 z^{-1})\allowbreak(1+\tfrac12 z^{-1})\bigr]$. Confirm $h[n]\to 0$ and identify
the slowest-decaying mode.
\end{exercise}

\begin{exercise}
Show that a real-coefficient filter must have poles and zeros that are either real or come
in complex-conjugate pairs. (Hint: take the complex conjugate of the polynomial equation
$A(p)=0$ and use that the coefficients are real.)
\end{exercise}

\begin{exercise}
An analog RC lowpass has $H(s)=1/(1+sRC)$. Find its pole in the s-plane, its cutoff
$\omega_c$, and the magnitude $\abs{H(\jj\omega)}$. Sketch the pole and explain why the
system is unconditionally stable.
\end{exercise}

\begin{exercise}
Using the sampling relation $z=e^{sT}$ with $T=1$ ms, map the analog pole $s=-200+\jj
2\pi(1000)$ into the z-plane. Compute $\abs{z}$ and the angle; is the resulting discrete
pole inside the unit circle? What discrete-time resonant frequency (in rad/sample) does it
correspond to?
\end{exercise}

\begin{exercise}
Place two zeros exactly on the unit circle at $z=e^{\pm\jj\pi/2}$ and no poles (other than
at the origin). Write $H(z)$, identify what frequency it nulls, and explain in one
sentence how you would use this as a notch to kill a single interfering tone in an EMG
recording (e.g. mains pickup after appropriate frequency scaling).
\end{exercise}

\begin{exercise}
\textbf{(Build/reflection.)} On an 8-bit microcontroller, suppose you implement
$y[n]=2r\cos(\omega_0)\,y[n-1]\allowbreak-r^2 y[n-2]+x[n]$ with $r=0.999$. Explain why this is risky
in fixed-point arithmetic, what failure you would observe, and how moving the pole inward
(smaller $r$) trades resonance sharpness for safety. Relate your answer to the wrist
controller's tap-detection filter.
\end{exercise}
